This thesis develops SNPbag, a transformer-based model for binary
phenotype prediction from genotype data, and compares it against
LDpred2-auto, a Bayesian polygenic risk score method. The model uses a
BERT-style encoder trained in two stages: first by reconstructing
randomly masked genotypes from their surrounding genomic context, then
by fine-tuning for phenotype prediction. Both methods are evaluated on
synthetic data from the HAPNEST project, comprising $100000$
individuals and $14000$ SNPs from chromosome~22, with phenotypes
simulated under a logistic disease model with $500$ causal variants and
$10\%$ prevalence.

Pre-training validation accuracy improved from $84.18\%$ at $1000$
individuals to $90.19\%$ at $100000$, with no signs of overfitting,
indicating that the model captures meaningful genomic structure and has
not yet saturated. For phenotype prediction, neither method achieved
strong discriminative performance: SNPbag obtained a test AUC of $0.54$
and LDpred2-auto $0.51$, both only marginally above the random baseline
of $0.5$. These results reflect the severe computational constraints
under which both methods operated — SNPbag was fine-tuned for a single
epoch with a frozen encoder, and LDpred2-auto was limited to $14000$
SNPs from a single chromosome — rather than a fundamental limitation of
either approach. Extending the model to longer sequences, more training
data, and full encoder fine-tuning are the most important directions for
future work.